
\par\noindent\textit{命题~\ref{prop:7-2-217} 的证明。}
其证明与定理~\ref{thm:7-2-212} 的证明类似。
为完整起见，仍写出细节。沿用相同的归纳方案，
证明如下引理。

\begin{lemma}{7.2.243}\label{lem:7-2-243}
在命题~\ref{prop:7-2-217} 的情形下，有：

\par\noindent\textup{(7.2.244.1)}
$\widetilde{\mathfrak C}'_i$
可提升为 $[0,1]^2\times C'_i$ 的
$A_{\bar n(\lambda_{i+\ell}),\bar k(\lambda_{i+\ell})}$
类模型。特别地，
$\operatorname{Incl}\colon C'_i\to
\widetilde{\mathfrak C}'_i$ 与
$\operatorname{Incl}_s\colon C'_i\to
(\widetilde{\mathfrak C}'_i)_{r=r_0}$
均可提升为滤过
$A_{\bar n(\lambda_{i+\ell}),\bar k(\lambda_{i+\ell})}$
同态。

\par\noindent\textup{(7.2.244.2)}
$\operatorname{Incl}_s\colon C'_i\to
\widehat{\mathfrak C}'_{i,r=0}$
可提升为滤过
$A_{\bar n(\lambda_{i+\ell}),\bar k(\lambda_{i+\ell})}$
同态。
$\widehat{\mathfrak C}'_{i,r=0}$
成为 $[0,1]_s\times C'_i$ 的
$A_{\bar n(\lambda_{i+\ell}),\bar k(\lambda_{i+\ell})}$
类模型。作为同一类的
$[0,1]_s\times C'_i$ 模型，
有
$(\widetilde{\mathfrak C}'_i)_{r=r_0}
=\widehat{\mathfrak C}'_{i,r=r_0}$。
\translatornote{原书此处及证明末段把任意
$r_0$ 对应的右端都写成
$\widehat{\mathfrak C}'_{i,r=0}$；
条件 (7.2.216.2) 及本引理的双端模型要求
右端随 $r_0$ 取相应端面。}

\par\noindent\textup{(7.2.244.3)}
$\operatorname{Incl}\colon C'_i\to
\widehat{\mathfrak C}'_i$
可提升为滤过
$A_{\bar n(\lambda_{i+\ell}),\bar k(\lambda_{i+\ell})}$
同态。
$\widehat{\mathfrak C}'_i$
成为 $[0,1]^2\times C'_i$ 的
$A_{\bar n(\lambda_{i+\ell}),\bar k(\lambda_{i+\ell})}$
类模型。
特别地，
$\widehat{\mathfrak C}'_{i,r=r_0}$
可提升为 $[0,1]_s\times C'_i$
的同阶模型。

此外有等式
\[
\operatorname{Incl}_s
=\operatorname{Eval}_{r=0}\circ
\operatorname{Incl},
\]
它是 (7.2.244.2) 中从
$C'_i$ 到
$\widehat{\mathfrak C}'_{i,r=0}$
的同态，并且

\par\noindent\textup{(7.2.244.4)}
$\operatorname{Eval}_{r=0}\circ
\widetilde{\mathfrak H}_i
=\mathcal I'_i\circ\operatorname{Incl}_s$
作为滤过
$A_{\bar n(\lambda_{i+\ell}),\bar k(\lambda_{i+\ell})}$
同态成立；

\par\noindent\textup{(7.2.244.5)}
$\operatorname{Eval}_{r=1}\circ
\widetilde{\mathfrak H}_i
=\operatorname{Incl}_s\circ I'_i$
作为同阶滤过同态成立；

\par\noindent\textup{(7.2.244.6)}
$\operatorname{Eval}_{s=s_0}\circ
\widetilde{\mathfrak H}_i
=\operatorname{Incl}_r\circ I'_i$
作为同阶滤过同态成立。
\end{lemma}

\begin{proof}
对 $(\ell,i)$ 作双重归纳。
假设引理对所有
$(\ell',i')<(\ell,i)$ 成立，现证明
$(\ell,i)$ 的情形。记
$(n',K')=(\bar n(\lambda_{i+\ell}),
\bar k(\lambda_{i+\ell}))$。
归纳假设给出某个 $(n,K)$，
使 (7.2.244) 在 $i$ 处作为
$A_{n,K}$ 结构成立\translatornote{原书此处写
“(30.244)”，但本段唯一的归纳陈述是
(7.2.244)。}；
在 $i+1$ 处，(7.2.244) 已作为
$A_{n',K'}$ 结构成立。
可假定 $((n,K),(n',K'))$ 属于
(7.2.130.1) 或 (7.2.130.2) 的类型。

取 $\beta\in\mathfrak G$、
$\lVert\beta\rVert=n'$，
考虑障碍类
\[
o_{K',\beta}(\widetilde{\mathfrak H}_i)
\in\operatorname{Hom}\bigl(
B_{K'}\overline C'_i[1],
\overline{\widetilde{\mathfrak C}}'_{i+1}[1]
\bigr).
\]
由于
$\widetilde{\mathfrak C}'_{i+1}$
在 $A_{n',K'}$ 层次满足 (7.2.244)，
故
\[
\begin{aligned}
\operatorname{Incl}_s\circ I'_i
&\colon C'_i\to
\widetilde{\mathfrak C}'_{i+1,r=1},\\
\operatorname{Incl}_r\circ I'_i
&\colon C'_i\to
\widetilde{\mathfrak C}'_{i+1,s=s_0}
\end{aligned}
\]
均为滤过 $A_{n',K'}$ 同态
（由引理~\ref{lem:7-2-184}，
$I'_i$ 和 $C'_i$ 已属 $A_\infty$ 类）。
由 (7.2.244.5)--(7.2.244.6)
在 $A_{n,K}$ 层次成立的归纳假设，得
\begin{equation}
\begin{aligned}
\overline{\operatorname{Eval}}_{r=1}\circ
o_{K',\beta}(\widetilde{\mathfrak H}_i)
+\delta_1\bigl(
(\operatorname{Incl}_s\circ I'_i)_{K',\beta}
\bigr)&=0,\\
\overline{\operatorname{Eval}}_{s=s_0}\circ
o_{K',\beta}(\widetilde{\mathfrak H}_i)
+\delta_1\bigl(
(\operatorname{Incl}_r\circ I'_i)_{K',\beta}
\bigr)&=0.
\end{aligned}
\tag{7.2.245}\label{eq:7-2-245}
\end{equation}
又注意到
\begin{equation}
\operatorname{Eval}_{s=s_0}\circ
\operatorname{Incl}_s\circ I'_i
=I'_i
=\operatorname{Eval}_{r=r_0}\circ
\operatorname{Incl}_r\circ I'_i.
\tag{7.2.246}\label{eq:7-2-246}
\end{equation}
置
\begin{equation}
\begin{aligned}
\mathfrak K
=\ker\Bigl(
&\overline{\widetilde{\mathfrak C}}'_{i+1,s=0}
\oplus
\overline{\widetilde{\mathfrak C}}'_{i+1,s=1}
\oplus
\overline{\widetilde{\mathfrak C}}'_{i+1,r=1}\\
&\xrightarrow{\ \overline B\ }
\bigoplus_{s_0=0}^{1}
\overline C'_{i+1,s=s_0,r=1}
\Bigr).
\end{aligned}
\tag{7.2.247}\label{eq:7-2-247}
\end{equation}
这里 $\overline B$ 由 (7.2.189.5)
中的 $B$ 诱导。
由子引理~\ref{sublem:7-2-228}，
同态
\begin{equation}
\bigl(
\overline{\operatorname{Eval}}_{s=0}
\oplus\overline{\operatorname{Eval}}_{s=1}
\oplus\overline{\operatorname{Eval}}_{r=1}
\bigr)
\colon
\overline{\widetilde{\mathfrak C}}'_{i+1}
\longrightarrow\mathfrak K
\tag{7.2.248}\label{eq:7-2-248}
\end{equation}
是满射的链同伦等价。
(7.2.246) 表明
\[
\bigl(
\operatorname{Incl}_r\circ I'_i,
\operatorname{Incl}_r\circ I'_i,
\operatorname{Incl}_s\circ I'_i
\bigr)
\in\operatorname{Hom}\bigl(
B_{K'}\overline C'_i[1],\mathfrak K
\bigr).
\]
因此，由 (7.2.245) 可取
$(\widetilde{\mathfrak H}_i)^I_{K',\beta}$
使
\begin{align}
\delta_1\bigl(
(\widetilde{\mathfrak H}_i)^I_{K',\beta}
\bigr)
+o_{K',\beta}(\widetilde{\mathfrak H}_i)
&=0,
\tag{7.2.249.1}\label{eq:7-2-249-1}\\
\overline{\operatorname{Eval}}_{r=1}
\bigl((\widetilde{\mathfrak H}_i)^I_{K',\beta}\bigr)
&=(\operatorname{Incl}_s\circ I'_i)_{K',\beta},
\tag{7.2.249.2}\label{eq:7-2-249-2}\\
\overline{\operatorname{Eval}}_{s=s_0}
\bigl((\widetilde{\mathfrak H}_i)^I_{K',\beta}\bigr)
&=(\operatorname{Incl}_r\circ I'_i)_{K',\beta}.
\tag{7.2.249.3}\label{eq:7-2-249-3}
\end{align}
(7.2.249.1) 允许用该分量将
$\widetilde{\mathfrak H}_i$
提升为滤过 $A_{n',K'}$ 同态，
记作 $\widetilde{\mathfrak H}_i^I$；
(7.2.249.2)--(7.2.249.3)
保证这个提升满足
(7.2.244.5)--(7.2.244.6)。

接下来提升滤过 $A_{n,K}$ 同态
\[
\operatorname{Incl}_s\colon C'_i\to
\widehat{\mathfrak C}'_{i,r=1},
\]
并修正 $\widetilde{\mathfrak H}_i^I$，
使 (7.2.244.4) 也作为滤过
$A_{n',K'}$ 同态等式成立。
考虑 $\operatorname{Incl}_s$ 的障碍类
\[
o_{K',\beta}(\operatorname{Incl}_s)
\in\operatorname{Hom}\bigl(
B_{K'}\overline C'_i[1],
\overline{\widehat{\mathfrak C}}'_{i,r=1}[1]
\bigr).
\]
由于下列等式已作为滤过 $A_{n,K}$
同态成立：
\begin{align}
\operatorname{Eval}_{r=0}\circ
\widetilde{\mathfrak H}_i^I
&=\mathcal I'_i\circ\operatorname{Incl}_s,
\tag{7.2.250.1}\label{eq:7-2-250-1}\\
\operatorname{Eval}_{s=s_0}\circ
\operatorname{Incl}_s
&=\mathrm{id},
\tag{7.2.250.2}\label{eq:7-2-250-2}
\end{align}
得到
\begin{equation}
o_{K',\beta}(\mathcal I'_i\circ
\operatorname{Incl}_s)
+\delta_1\Bigl(
\overline{\operatorname{Eval}}_{r=0}\circ
(\widetilde{\mathfrak H}_i)^I_{K',\beta}
\Bigr)=0.
\tag{7.2.250.3}\label{eq:7-2-250-3}
\end{equation}
这里 $\mathcal I'_i\circ
\operatorname{Incl}_s$
按滤过 $A_{n,K}$ 同态复合理解。

由于 $\mathcal I'_i$ 是滤过
$A_\infty$ 同态，而
$\operatorname{Incl}_s$ 是滤过
$A_{n,K}$ 同态，故
\begin{equation}
\begin{aligned}
o_{K',\beta}(\mathcal I'_i\circ
\operatorname{Incl}_s)
={}&-\delta_1\bigl((\mathcal I'_i)_{K',\beta}\bigr)
\circ
(\overline{\operatorname{Incl}}_s)_{1,\beta_0}^{\otimes K'}\\
&+(\overline{\mathcal I'_i})_1\circ
o_{K',\beta}(\operatorname{Incl}_s).
\end{aligned}
\tag{7.2.251}\label{eq:7-2-251}
\end{equation}
置
\[
E=-(\mathcal I'_i)_{K',\beta}\circ
(\overline{\operatorname{Incl}}_s)_{1,\beta_0}^{\otimes K'}
+\overline{\operatorname{Eval}}_{r=0}\circ
(\widetilde{\mathfrak H}_i)^I_{K',\beta}.
\]
由 (7.2.250.3)、(7.2.251) 得
\begin{equation}
(\overline{\mathcal I'_i})_1\circ
o_{K',\beta}(\operatorname{Incl}_s)
+\delta_1(E)=0.
\tag{7.2.252}\label{eq:7-2-252}
\end{equation}
再计算
\[
\begin{aligned}
\operatorname{Eval}_{s=s_0}\circ
\operatorname{Eval}_{r=0}\circ
\widetilde{\mathfrak H}_i^I
&=\operatorname{Eval}_{r=0}\circ
\operatorname{Eval}_{s=s_0}\circ
\widetilde{\mathfrak H}_i^I\\
&=\operatorname{Eval}_{r=0}\circ
\operatorname{Incl}_r\circ I'_i
=I'_i.
\end{aligned}
\]
这作为滤过 $A_{n',K'}$ 同态成立；
第二个等号所用的 (7.2.244.6)，
已在定义 $(\widetilde{\mathfrak H}_i)^I_{K',\beta}$
时提升至该层次。因此
\[
\begin{aligned}
\overline{\operatorname{Eval}}_{s=s_0}\circ
\overline{\operatorname{Eval}}_{r=0}\circ
(\widetilde{\mathfrak H}_i^I)_{K',\beta}
&=(I'_i)_{K',\beta}\\
&=(I'_i)_{K',\beta}\circ
\bigl(
\overline{\operatorname{Eval}}_{s=s_0}\circ
(\overline{\operatorname{Incl}}_s)_{1,\beta_0}
\bigr)^{\otimes K'}\\
&=\overline{\operatorname{Eval}}_{s=s_0}\circ
(\mathcal I'_i)_{K',\beta}\circ
(\overline{\operatorname{Incl}}_s)_{1,\beta_0}^{\otimes K'}.
\end{aligned}
\]
末个等号用到 (7.2.206)。
对 $E$ 的定义施以
$\overline{\operatorname{Eval}}_{s=s_0}$，
得到
\begin{equation}
\overline{\operatorname{Eval}}_{s=s_0}(E)=0.
\tag{7.2.253}\label{eq:7-2-253}
\end{equation}
另一方面，由
$\operatorname{Eval}_{s=s_0}\circ
\operatorname{Incl}_s=\mathrm{id}$
在 $A_{n,K}$ 同态意义下成立，
有
\begin{equation}
\overline{\operatorname{Eval}}_{s=s_0}
\bigl(o_{K',\beta}(\operatorname{Incl}_s)\bigr)=0.
\tag{7.2.254}\label{eq:7-2-254}
\end{equation}
注意 $(\overline{\mathcal I'_i})_1$
诱导链同伦等价
\begin{equation}
\begin{aligned}
&\bigl\{x\in
\overline{\widehat{\mathfrak C}}'_{i,r=1}
\mid\overline{\operatorname{Eval}}_{s=s_0}(x)=0,
\ s_0=0,1\bigr\}\\
&\quad\longrightarrow
\bigl\{x\in
\overline{\widehat{\mathfrak C}}'_{i+1,r=0}
\mid\overline{\operatorname{Eval}}_{s=s_0}(x)=0,
\ s_0=0,1\bigr\}.
\end{aligned}
\tag{7.2.255}\label{eq:7-2-255}
\end{equation}
由 (7.2.252)--(7.2.255)，
可取
\[
(\operatorname{Incl}_s)^I_{K',\beta}
\in\operatorname{Hom}\bigl(
B_{K'}\overline C'_i[1],
\overline{\widehat{\mathfrak C}}'_{i,r=1}[1]
\bigr)
\]
满足
\begin{align}
o_{K',\beta}(\operatorname{Incl}_s)
+\delta_1\bigl(
(\operatorname{Incl}_s)^I_{K',\beta}
\bigr)&=0,
\tag{7.2.256.1}\label{eq:7-2-256-1}\\
\overline{\operatorname{Eval}}_{s=s_0}
\bigl((\operatorname{Incl}_s)^I_{K',\beta}\bigr)
&=0.
\tag{7.2.256.2}\label{eq:7-2-256-2}
\end{align}
由 (7.2.256.1)，
用该分量将 $\operatorname{Incl}_s$
提升为滤过 $A_{n',K'}$ 同态，
记作 $\operatorname{Incl}_s^I$；
(7.2.256.2) 保证
\[
\operatorname{Eval}_{s=s_0}\circ
\operatorname{Incl}_s^I=\mathrm{id}
\]
也在该层次成立。

接下来利用 (7.2.244.4)
在 $A_{n,K}$ 层次成立这一事实。
$\operatorname{Eval}_{r=0}\circ
\widetilde{\mathfrak H}_i^I
-\mathcal I'_i\circ
\operatorname{Incl}_s^I$
在 $B_{K'}\overline C'_i[1]$ 上的限制，
模去 $T^{\omega\cap\beta}
\Lambda^+_{0,\mathrm{nov}}$ 后
可写为
\[
\sum_{\substack{\beta\in\mathfrak G\\
\lVert\beta\rVert=n'}}
T^{\beta\cap\omega}e^{\mu(\beta)/2}
\bigl(
\operatorname{Eval}_{r=0}\circ
\widetilde{\mathfrak H}_i^I
-\mathcal I'_i\circ
\operatorname{Incl}_s^I
\bigr)_{K',\beta},
\]
其中
\[
\bigl(
\operatorname{Eval}_{r=0}\circ
\widetilde{\mathfrak H}_i^I
-\mathcal I'_i\circ
\operatorname{Incl}_s^I
\bigr)_{K',\beta}
\in\operatorname{Hom}\bigl(
B_{K'}\overline C'_i[1],
\overline{\widehat{\mathfrak C}}'_{i+1,r=0}[1]
\bigr).
\]
由 (7.2.244.4) 在 $A_{n,K}$ 层次成立，
得到
\begin{equation}
\delta_1\Bigl(
\bigl(\operatorname{Eval}_{r=0}\circ
\widetilde{\mathfrak H}_i^I
-\mathcal I'_i\circ
\operatorname{Incl}_s^I\bigr)_{K',\beta}
\Bigr)=0.
\tag{7.2.257}\label{eq:7-2-257}
\end{equation}
又由
\[
\operatorname{Eval}_{s=s_0}\circ
\operatorname{Eval}_{r=0}\circ
\widetilde{\mathfrak H}_i^I
=I'_i
=\operatorname{Eval}_{s=s_0}\circ
\mathcal I'_i\circ
\operatorname{Incl}_s^I
\]
作为滤过 $A_{n',K'}$ 同态成立
（末个等号用到 (7.2.206)），得
\begin{equation}
\overline{\operatorname{Eval}}_{s=s_0}
\Bigl(
\bigl(\operatorname{Eval}_{r=0}\circ
\widetilde{\mathfrak H}_i^I
-\mathcal I'_i\circ
\operatorname{Incl}_s^I\bigr)_{K',\beta}
\Bigr)=0.
\tag{7.2.258}\label{eq:7-2-258}
\end{equation}
由 (7.2.257)--(7.2.258)
及 (7.2.255) 的链同伦等价性，
可取
\[
\begin{aligned}
\operatorname{Corr1}
&\in\operatorname{Hom}\bigl(
B_{K'}\overline C'_i[1],
\overline{\widehat{\mathfrak C}}'_{i,r=1}[1]
\bigr),\\
\operatorname{Corr2}
&\in\operatorname{Hom}\bigl(
B_{K'}\overline C'_i[1],
\overline{\widehat{\mathfrak C}}'_{i+1,r=0}[1]
\bigr),
\end{aligned}
\]
满足
\begin{align}
\delta_1(\operatorname{Corr1})&=0,
\tag{7.2.259.1}\label{eq:7-2-259-1}\\
\overline{\operatorname{Eval}}_{r=0}\circ
(\widetilde{\mathfrak H}_i^I)_{K',\beta}
-(\mathcal I'_i\circ
\operatorname{Incl}_s^I)_{K',\beta}
&=(\overline{\mathcal I'_i})_1\circ
\operatorname{Corr1}
-\delta_1(\operatorname{Corr2}),
\tag{7.2.259.2}\label{eq:7-2-259-2}\\
\overline{\operatorname{Eval}}_{s=s_0}\circ
\operatorname{Corr1}
&=\overline{\operatorname{Eval}}_{s=s_0}\circ
\operatorname{Corr2}=0.
\tag{7.2.259.3}\label{eq:7-2-259-3}
\end{align}
由图式~\ref{diag:7-2-6} 的正合性，
可取
\[
\operatorname{Corr3}\in
\operatorname{Hom}\bigl(
B_{K'}\overline C'_i[1],
\overline{\widetilde{\mathfrak C}}'_i[1]
\bigr)
\]
使
\[
\begin{aligned}
\overline{\operatorname{Eval}}_{r=0}\circ
\operatorname{Corr3}
&=\operatorname{Corr2},\\
\overline{\operatorname{Eval}}_{r=1}\circ
\operatorname{Corr3}
&=0,\qquad
\overline{\operatorname{Eval}}_{s=s_0}\circ
\operatorname{Corr3}=0.
\end{aligned}
\]
现在置
\[
\begin{aligned}
(\widetilde{\mathfrak H}_i)_{K',\beta}
&=(\widetilde{\mathfrak H}_i^I)_{K',\beta}
+\delta_1(\operatorname{Corr3}),\\
(\operatorname{Incl}_s)_{K',\beta}
&=(\operatorname{Incl}_s^I)_{K',\beta}
+\operatorname{Corr1}.
\end{aligned}
\]
直接核对可知二者满足 (7.2.244.4)。
至此，已把
$\widetilde{\mathfrak H}_i\colon
C'_i\to\widetilde{\mathfrak C}'_{i+1}$
及
$\operatorname{Incl}_s\colon C'_i\to
\widehat{\mathfrak C}'_{i,r=1}$
提升，使 (7.2.244.4)--(7.2.244.6)
在 $A_{n',K'}$ 层次成立。

下面提升两个映射
\begin{align}
\operatorname{Incl}\colon
C'_i&\longrightarrow
\widehat{\mathfrak C}'_i,
\tag{7.2.260.1}\label{eq:7-2-260-1}\\
\operatorname{Incl}_s\colon
C'_i&\longrightarrow
\widehat{\mathfrak C}'_{i,r=0},
\tag{7.2.260.2}\label{eq:7-2-260-2}
\end{align}
使它们连同
$\operatorname{Incl}_s\colon C'_i\to
\widehat{\mathfrak C}'_{i,r=1}$
令滤过 $A_{n',K'}$ 代数
$\widehat{\mathfrak C}'_i$
成为 $[0,1]^2\times C'_i$ 的同阶模型。

为构造 (7.2.260)，回顾引理~\ref{lem:7-2-208}
中已定义滤过 $A_\infty$ 同态
\[
\operatorname{Incl}_r\colon C'_i\to
\widehat{\mathfrak C}'_{i,s=s_0}.
\]
又有
\[
\operatorname{Eval}_{r=1}\circ
\operatorname{Incl}_r
=\mathrm{id}
=\operatorname{Eval}_{s=s_0}\circ
\operatorname{Incl}_s
\colon C'_i\to C'_i.
\]
因此可由子引理~\ref{sublem:7-2-228}
把 (7.2.260.1) 中的
$\operatorname{Incl}$
提升为滤过 $A_{n',K'}$ 同态，
并使
\[
\operatorname{Eval}_{r=1}\circ
\operatorname{Incl}
=\operatorname{Incl}_s,
\qquad
\operatorname{Eval}_{s=s_0}\circ
\operatorname{Incl}
=\operatorname{Incl}_r
\]
作为同阶滤过同态成立。
随后定义 (7.2.260.2) 中的映射为
\[
\operatorname{Incl}_s
=\operatorname{Eval}_{r=0}\circ
\operatorname{Incl}.
\]
这就定义了 (7.2.260) 的两个同态。

为完成引理~\ref{lem:7-2-243} 的证明，
只需再提升
$\widetilde{\mathfrak C}'_i$
及
$\operatorname{Incl}\colon C'_i\to
\widetilde{\mathfrak C}'_i$，
使 (7.2.244.1)--(7.2.244.2) 成立。

\begin{statement}{子引理}{7.2.261}\label{sublem:7-2-261}\itshape
可将
$\widetilde{\mathfrak C}'_{i,s=s_0}$
上的滤过 $A_{n,K}$ 结构提升为
滤过 $A_{n',K'}$ 结构，
使
$\operatorname{Eval}_{r=r_0}\colon
\widetilde{\mathfrak C}'_{i,s=s_0}
\to C'_i$
为滤过 $A_{n',K'}$ 同态。
还可提升
$\operatorname{Incl}_r\colon C'_i\to
\widetilde{\mathfrak C}'_{i,s=s_0}$
为滤过 $A_{n',K'}$ 同态，
使
$\bigl(\widetilde{\mathfrak C}'_{i,s=s_0};
\operatorname{Eval}_{r=r_0},
\operatorname{Incl}_r\bigr)$
成为 $[0,1]_r\times C'_i$ 的
$A_{n',K'}$ 类模型。
\end{statement}

子引理~\ref{sublem:7-2-261}
是命题~\ref{prop:7-2-190} 证明中
归纳步骤的单参数版本；
命题~\ref{prop:7-2-190} 则是双参数版本。
单参数版本与双参数版本类似而较易，
故略去其证明。

注意
$\widetilde{\mathfrak C}'_{i,r=r_0}
=\widehat{\mathfrak C}'_{i,r=r_0}$
已是滤过 $A_\infty$ 代数。
在 $r_0=0$ 端面，由 (7.2.260.2)
已把相应的 $\operatorname{Incl}_s$
提升为滤过 $A_{n',K'}$ 同态；
$r_0=1$ 端面的提升已在前文完成。

于是，命题~\ref{prop:7-2-190}
证明中的归纳步骤表明：
可将
$\widetilde{\mathfrak C}'_i$
及
$\operatorname{Incl}\colon C'_i\to
\widetilde{\mathfrak C}'_i$
分别提升为滤过 $A_{n',K'}$ 结构
与滤过 $A_{n',K'}$ 同态，
使 (7.2.244.1)--(7.2.244.2) 成立。
引理~\ref{lem:7-2-243} 得证。
\end{proof}

命题~\ref{prop:7-2-217}
由引理~\ref{lem:7-2-243} 立即得出。

\begin{remark}{7.2.262}\label{rem:7-2-262}
本小节的讨论可以推广，
以处理同伦的同伦的同伦等情形。
为完整构造弗洛尔上同调的族版本，
似乎需要这种推广（见\textbf{[Fuk02I]}）。
为此似乎有必要进一步发展一般理论，
使本小节的内容更加系统。
本书所发展的（滤过）
$A_\infty$ 代数同伦理论，
正是为适应这种推广而设计的。
\end{remark}
